\documentclass[12pt,a4paper]{article}
 
\usepackage[margin=2.5cm,headheight=15pt]{geometry}
\usepackage{graphicx}
\usepackage{float}
\usepackage{hyperref}
\usepackage{xcolor}
\usepackage{listings}
\usepackage{titlesec}
\usepackage{booktabs}
\usepackage{enumitem}
\usepackage{fancyhdr}
\usepackage{setspace}
\usepackage{parskip}
\usepackage{caption}
\usepackage{subcaption}
\usepackage{mdframed}
\usepackage{array}
\usepackage{longtable}
 

\definecolor{codebg}{RGB}{245,245,245}
\definecolor{codeframe}{RGB}{200,200,200}
\definecolor{keywordcolor}{RGB}{0,0,180}
\definecolor{commentcolor}{RGB}{0,128,0}
\definecolor{stringcolor}{RGB}{163,21,21}
\definecolor{sectioncolor}{RGB}{30,30,120}
\definecolor{accentblue}{RGB}{0,70,180}
 

\titleformat{\section}{\large\bfseries\color{sectioncolor}}{}{0em}{\thesection\quad}[\titlerule]
\titleformat{\subsection}{\normalsize\bfseries\color{accentblue}}{}{0em}{\thesubsection\quad}
\titlespacing*{\section}{0pt}{16pt}{6pt}
\titlespacing*{\subsection}{0pt}{10pt}{4pt}

\lstset{
  backgroundcolor=\color{codebg},
  frame=single,
  rulecolor=\color{codeframe},
  basicstyle=\ttfamily\footnotesize,
  keywordstyle=\color{keywordcolor}\bfseries,
  commentstyle=\color{commentcolor}\itshape,
  stringstyle=\color{stringcolor},
  breaklines=true,
  showstringspaces=false,
  tabsize=4,
  numbers=left,
  numberstyle=\tiny\color{gray},
  numbersep=6pt,
  xleftmargin=12pt
}

\pagestyle{fancy}
\fancyhf{}
\rhead{\small Mini Game Hub}
\lhead{\small SSL Project Report}
\rfoot{\small Page \thepage}
\lfoot{\small 2025--26}
\renewcommand{\headrulewidth}{0.4pt}
\renewcommand{\footrulewidth}{0.4pt}
 
\hypersetup{
  colorlinks=true,
  linkcolor=accentblue,
  urlcolor=accentblue,
  citecolor=accentblue
}
 

\begin{document}
 

\begin{titlepage}
  \centering
  \vspace*{1.5cm}
 
  \includegraphics[width=0.85\textwidth]{report/gamehub.jpeg}
 
  \vspace{1.2cm}
  {\Huge\bfseries\color{sectioncolor} Mini Game Hub}\\[0.4cm]
  {\large\bfseries Project Report}\\[0.2cm]
  {\normalsize SSL Project}
 
  \vfill
  \begin{tabular}{rl}
    \textbf{Authors:}   & Keshav \quad Arnav \\
      
    \textbf{Language:}  & Python , Bash      \\
  \end{tabular}
  \vspace{1cm}
\end{titlepage}
 
\tableofcontents
\newpage

\section{External Libraries Used}

 
The project relies on the following external libraries and tools. All Python
dependencies can be installed via \texttt{pip} and are documented in
\texttt{requirements.txt}.
 
\begin{longtable}{>{\bfseries}p{3.2cm} p{2.2cm} p{8cm}}
  \toprule
  Library & Version & Purpose \\
  \midrule
  \endhead
  
  pygame-ce & 2.5.7 & Core game framework. Handles the window, event loop, image
                 rendering, fonts, and all user input across every screen. \\[4pt]
  matplotlib & 3.10.8 & Generates the bar chart (player rankings) and pie chart
                      (game distribution) saved as PNG files that pygame then
                      displays. \\[4pt]
  numpy & 2.4.4 & Board representation for all three games using 2-D integer
                arrays. Used extensively for win-condition detection via
                \texttt{sliding\_window\_view} and \texttt{np.diagonal}. \\[4pt]
  
  subprocess & (stdlib) & Runs the Bash leaderboard script
                           (\texttt{leaderboard.sh}) from within Python,
                           capturing terminal output. \\[4pt]
  datetime & (stdlib) & Provides today's date string written into
                         \texttt{history.csv} when a result is recorded. \\[4pt]
  sys & (stdlib) & Reads command-line arguments (player names) passed from
                    \texttt{main.sh}. \\[4pt]
  sha256sum & (Bash) & Hashes passwords before storing them in
                        \texttt{users.tsv}, providing basic security without
                        storing plain text. \\
  \bottomrule
  \caption{External libraries and their roles.}
\end{longtable}
 

\section{Design and Architecture}

 
\subsection{Overall Structure}
 
The project is split into a Bash entry-point layer and a Python/pygame
application layer. The separation means that authentication and password
management happen entirely in the shell before Python is ever invoked.
 
\begin{figure}[H]
  \centering
  \includegraphics[width=0.82\textwidth]{report/arcade_bg.jpeg}
  \caption{Retro arcade aesthetic used throughout the project.}
\end{figure}
 
\subsection{BaseGame Class}
 
A single abstract base class, \texttt{BaseGame} (in \texttt{game.py}), defines
the interface shared by all three games:
 
\begin{lstlisting}[language=Python, caption={BaseGame skeleton}]
class BaseGame:
    def __init__(self, player1, player2):
        self.player1 = player1
        self.player2 = player2
        self.current_player = 1
        self.board = None
        self.running = True
        self.game_over = False
 
    def switch_turn(self):
        self.current_player = 2 if self.current_player == 1 else 1
 
    def check_win(self):  pass
    def reset(self):      pass
\end{lstlisting}
 
Each game subclass inherits the two-player state and the turn-switching helper,
then overrides \texttt{check\_win}, \texttt{reset}, and adds its own
\texttt{run(screen)} method which contains the pygame event loop.
 
\subsection{Data Flow}
 
\begin{enumerate}[leftmargin=*, label=\arabic*.]
  \item User invokes \texttt{bash main.sh}.
  \item Both players authenticate via \texttt{login()} / \texttt{register()}.
  \item \texttt{python3 game.py \$username1 \$username2} is launched.
  \item The Game Hub menu appears; the user picks a game.
  \item After each game, \texttt{record\_result()} appends a CSV row to
        \texttt{history.csv}.
  \item \texttt{show\_sort\_menu()} asks the user how to sort the leaderboard.
  \item \texttt{leaderboard.sh} is executed via \texttt{subprocess.run()} and
        prints to the terminal.
  \item \texttt{generate\_graphs()} produces two PNG files; \texttt{show\_graphs()}
        displays them inside the pygame window.
\end{enumerate}
 

\section{Features Implemented}

 
\subsection{Authentication System}
 
\begin{itemize}[leftmargin=*]
  \item \textbf{Registration}: New users choose a username and password. The
        password is hashed with \texttt{sha256sum} and stored as a tab-separated
        record in \texttt{users.tsv}. Plain-text passwords are never written
        to disk.
  \item \textbf{Login}: Existing usernames are looked up in \texttt{users.tsv}.
        The supplied password is hashed and compared; the loop repeats until
        correct.
  \item \textbf{Input validation}: Usernames are validated with a regex
        \texttt{[a-zA-Z0-9]+\$} before any lookup occurs.
  \item \textbf{Duplicate prevention}: The shell rejects two players using the
        same username.
\end{itemize}
 
\subsection{Game Hub Menu}
 
A custom pygame screen (\texttt{game\_hub.png} background) presents three
game buttons. Each button scales up by 10\% on hover, implemented by pre-computing
a smooth-scaled surface and swapping it in when the mouse rectangle collides.
 
\subsection{Connect 4}
 
\begin{figure}[H]
  \centering
  \includegraphics[width=0.78\textwidth]{report/connect4.jpeg}
  \caption{Connect 4 game screen with the 7\texttimes7 grid.}
\end{figure}
 
\begin{itemize}[leftmargin=*]
  \item 7$\times$7 board stored in a \texttt{numpy} integer array.
  \item Gravity: \texttt{get\_next\_open\_row()} scans upward from the bottom to
        find the first empty cell in a column.
  \item Win detection uses \texttt{np.lib.stride\_tricks.sliding\_window\_view}
        for all four directions (horizontal, vertical, both diagonals) in a
        vectorised, loop-free manner.
  \item Draw detection: \texttt{np.any(self.board == 0)} — if no empty cell
        remains the game is a draw.
  \item Animated turn indicator and large win/draw overlay text.
\end{itemize}
 
\subsection{Tic-Tac-Toe (10\texttimes10)}
 
\begin{figure}[H]
  \centering
  \includegraphics[width=0.78\textwidth]{report/tictactoe.jpeg}
  \caption{Tic-Tac-Toe on a 10\texttimes10 grid, win requires 5 in a row.}
\end{figure}
 
\begin{itemize}[leftmargin=*]
  \item Extended variant: 10$\times$10 grid, five-in-a-row required to win.
  \item Pieces drawn as cyan X shapes and pink circles with pixel-perfect
        padding computed from the cell size.
  \item A winning line is drawn across the five matching cells.
  \item Press \texttt{R} to reset the board without leaving the screen.
\end{itemize}
 
\subsection{Othello (Reversi)}
 
\begin{figure}[H]
  \centering
  \begin{subfigure}[b]{0.48\textwidth}
    \includegraphics[width=\textwidth]{report/othello.jpeg}
    \caption{Othello board at start of a game.}
  \end{subfigure}
  \hfill
  \begin{subfigure}[b]{0.48\textwidth}
    \includegraphics[width=\textwidth]{report/gamehub.jpeg}
    \caption{Game Hub menu screen.}
  \end{subfigure}
  \caption{Othello and Game Hub screenshots.}
\end{figure}
 
\begin{itemize}[leftmargin=*]
  \item Standard 8$\times$8 Othello rules with full disc-flipping logic.
  \item Valid moves are computed each turn by scanning all eight directions
        from every empty cell for an opponent sandwich.
  \item If a player has no valid moves their turn is skipped; if neither
        player can move the game ends and disc counts determine the winner.
  \item Grey dot hints show valid placement positions.
  \item Scores (disc counts) displayed above the board alongside player names.
\end{itemize}
 
\subsection{Result Recording}
 
\begin{lstlisting}[language=Python, caption={Recording a game result to CSV}]
def record_result(winner, loser, game_name):
    if winner is None:
        return
    with open("history.csv", "a") as f:
        f.write(f"{winner},{loser},{date.today().isoformat()},{game_name}\n")
\end{lstlisting}
 
Every completed game appends one line:
\texttt{winner,loser,YYYY-MM-DD,GameName}.
Draws and quit events return special sentinel values and are filtered out before
writing.
 
\subsection{Leaderboard}
 
\begin{itemize}[leftmargin=*]
  \item A Bash AWK script (\texttt{leaderboard.sh}) reads \texttt{history.csv},
        aggregates wins, losses, and W/L ratio per player per game, and prints
        a formatted table to the terminal.
  \item The user selects a sort metric (wins / losses / ratio) from a dedicated
        pygame sort-menu screen before the script runs.
  \item The same metric is passed to \texttt{generate\_graphs()} so the bar
        chart reflects the chosen ranking.
\end{itemize}
 
\subsection{Graphical Leaderboard}
 
\begin{figure}[H]
  \centering
  \includegraphics[width=0.72\textwidth]{report/leaderboard.jpeg}
  \caption{Wins Leaderboard bar chart rendered with matplotlib.}
\end{figure}
 
\begin{itemize}[leftmargin=*]
  \item \textbf{Bar chart}: Player totals for the chosen metric, rendered with
        the \texttt{PressStart2P} pixel font, transparent background, yellow
        labels, and coloured bars (green / red / cyan per metric).
  \item \textbf{Pie chart}: Distribution of games played across Connect 4,
        Tic-Tac-Toe, and Othello.
  \item Both PNGs are loaded into pygame, scaled, and drawn over a semi-transparent
        black panel so they are readable against any background.
\end{itemize}
 
\section{Code Logic Explained}

 
\subsection{BaseGame and Inheritance}
 
\texttt{BaseGame} stores the minimum shared state: two player names, whose turn
it is, and the game-over flag. Subclasses call \texttt{super().\_\_init\_\_()} to
initialise this state, then add their own board, fonts, and background image.
The \texttt{switch\_turn()} helper flips \texttt{current\_player} between 1 and
2 with a single ternary expression, avoiding any state inconsistency.
 
\subsection{Connect 4 Win Detection}
 
\begin{lstlisting}[language=Python, caption={Vectorised four-direction win check}]
def check_win(self, piece):
    # Horizontal
    w = np.lib.stride_tricks.sliding_window_view(self.board, (1, 4))
    if np.any(np.all(w == piece, axis=3)): return True
    # Vertical
    w = np.lib.stride_tricks.sliding_window_view(self.board, (4, 1))
    if np.any(np.all(w == piece, axis=2)): return True
    # Diagonal top-left to bottom-right
    w = np.lib.stride_tricks.sliding_window_view(self.board, (4, 4))
    if np.any(np.all(np.diagonal(w, axis1=2, axis2=3) == piece, axis=2)):
        return True
    # Diagonal top-right to bottom-left
    w = np.lib.stride_tricks.sliding_window_view(np.fliplr(self.board), (4,4))
    if np.any(np.all(np.diagonal(w, axis1=2, axis2=3) == piece, axis=2)):
        return True
    return False
\end{lstlisting}
 
\texttt{sliding\_window\_view} creates a view of every possible 1$\times$4 (or
4$\times$1 / 4$\times$4) sub-array without copying data.
\texttt{np.all(..., axis=N)} checks whether all four values equal \texttt{piece}.
The anti-diagonal case is handled by horizontally flipping the board before
extracting the main diagonal.
 
\subsection{Tic-Tac-Toe Win Detection}
 
The same stride-trick approach is used with window size 5 for the 10$\times$10
board. The function also records the start and end cell of the winning window so
that \texttt{draw\_win\_line()} can draw a line across it.
 
\subsection{Othello Flip Logic}
 
\begin{lstlisting}[language=Python, caption={Disc flipping in Othello}]
def flip_discs(self, r, c, player):
    opponent = 2 if player == 1 else 1
    for dr, dc in DIRS:           # 8 directions
        path = []
        nr, nc = r + dr, c + dc
        while inside(nr, nc) and board[nr, nc] == opponent:
            path.append((nr, nc))
            nr += dr; nc += dc
        if inside(nr, nc) and board[nr, nc] == player:
            for pr, pc in path:
                board[pr, pc] = player
\end{lstlisting}
 
For each of the eight compass directions the code walks outward collecting
opponent pieces. When it reaches a piece belonging to the current player the
entire collected path is flipped. If the walk hits the board edge or an empty
cell first, nothing is flipped in that direction.
 
\subsection{Authentication (Bash)}
 
\begin{lstlisting}[language=bash, caption={Password hashing in main.sh}]
passwordhash=$(echo -n "$password" | sha256sum | awk '{print $1}')
echo -e "${username}\t${passwordhash}" >> users.tsv
\end{lstlisting}
 
The password is piped through \texttt{sha256sum}; \texttt{awk} strips the
trailing filename that \texttt{sha256sum} appends. On login the same hash is
computed from user input and compared with \texttt{grep} + \texttt{cut}.
 
\subsection{Leaderboard Shell Script}
 
The AWK program in \texttt{leaderboard.sh} maintains four associative arrays
(\texttt{wins}, \texttt{losses}, \texttt{seen}) keyed by
\texttt{player|game}.
In the \texttt{END} block it iterates over \texttt{seen}, computes the W/L
ratio (using 9999 as a sentinel for infinity when losses are zero), and prints
each row. The output is piped to \texttt{sort -k\${SORT\_COL}rn} so that the
chosen metric column is sorted descending, then formatted by a second
\texttt{awk} pass.
 
\subsection{Graph Generation}
 
\texttt{generate\_graphs(sort\_by)} reads the same CSV, builds parallel lists of
players, wins, and losses, then creates two matplotlib figures with transparent
backgrounds. Bar values are annotated above each bar using \texttt{ax.text()}.
Both figures use \texttt{bbox\_inches='tight'} when saving so that long tick
labels are not clipped. The resulting PNGs are loaded by pygame in
\texttt{show\_graphs()}, placed over semi-transparent black panels for contrast,
and displayed until the user clicks anywhere.
 

\section{Directory Structure}
\begin{mdframed}
\begin{footnotesize}
\begin{verbatim}
.
|-- README.md
|-- main.sh                  Entry point: authenticate two players, launch game
|-- game.py                  BaseGame class, hub menu, leaderboard functions
|-- leaderboard.sh           Bash/AWK leaderboard aggregator and printer
|-- history.csv              Append-only match log (winner,loser,date,game)
|-- users.tsv                Username TAB sha256 password hash store
|-- requirements.txt         Python pip dependencies
|-- PressStart2P-Regular.ttf Pixel font for pygame and matplotlib
|-- games/
|   |-- connect4.py          Connect 4 game logic and rendering
|-- tictactoe.py         10x10 Tic-Tac-Toe logic and rendering
|    |-- othello.py           Othello (Reversi) logic and rendering
|    |-- __init__.py
|    |-- button_images
|        |-- connect4_button_cropped.png   connect4 button
|        |-- losses_button_transparent.png    losses button
|        |-- othello_button_transparent.png     othello button
|        |-- tictactoe_button_transparent.png    tictactoe button
|        |-- wins_button_transparent.png      wins button
|        |-- wlratio_button_transparent.png   wlratio button
|    |-- bg_images/
|        |-- tictactoe.png    Tic-Tac-Toe background
|        |-- othello.png      Othello background
|        |-- Connect4_final.png     connect4 background
|        |-- game_hub.png         mainmenu background
|        |-- genbg_resized.jpeg    graphs background
|        |-- sortmetric_bg.png      sortmetric menu background
|-- Report
     |-- report.tex
     |-- arcade_bg.jpeg
     |-- connect4.jpeg
     |-- gamehub.jpeg
     |-- leaderboard.jpeg
     |--tictactoe.jpeg
     |-- othello.jpeg
\end{verbatim}
\end{footnotesize}
\end{mdframed}
 

\section{Installation and Setup}

 
\subsection{Dependencies}
 
\begin{lstlisting}[language=bash, caption={Installing Python dependencies}]
pip install pygame matplotlib numpy
\end{lstlisting}
 
A \texttt{requirements.txt} is provided at the project root:
 
\begin{lstlisting}[caption={requirements.txt}]
pygame>=2.1.0
matplotlib>=3.5.0
numpy>=1.22.0
\end{lstlisting}
 
\subsection{System Requirements}
 
\begin{itemize}[leftmargin=*]
  \item Python 3.9 or later
  \item Bash 4.x (standard on Linux/macOS; use WSL on Windows)
  \item \texttt{sha256sum} (GNU coreutils, pre-installed on Linux/macOS)
  \item \texttt{awk} (gawk or mawk)
  \item A display server (X11 / Wayland / macOS Quartz) for pygame
\end{itemize}
 
\subsection{Running the Project}
 
\begin{lstlisting}[language=bash, caption={Launching the game}]
chmod +x main.sh leaderboard.sh
bash main.sh
\end{lstlisting}
 
The script prompts for two usernames and passwords, then opens the pygame
window. All game assets must be in the same directory as \texttt{game.py} (or
in \texttt{games/} for per-game assets).
 

 

\section{Usage}

 
\begin{enumerate}[leftmargin=*, label=\textbf{Step \arabic*.}]
  \item Run \texttt{bash main.sh} in a terminal.
  \item Enter a username for Player 1. If new, you will be prompted to create
        a password; if existing, enter your password.
  \item Repeat for Player 2 (must be a different username).
  \item The \textbf{Game Hub} window opens. Click on Connect 4, Tic-Tac-Toe,
        or Othello to start a game.
  \item Play the game. When it ends, a win/draw banner appears; click anywhere
        to continue.
  \item The \textbf{Sort Metric} screen appears. Click \emph{Wins},
        \emph{Losses}, or \emph{W/L Ratio} to choose how the leaderboard
        is sorted.
  \item The leaderboard is printed to the terminal and a graphical view
        (bar chart + pie chart) is shown in the pygame window.
  \item Click anywhere on the leaderboard screen to return to the Game Hub
        and play again.
\end{enumerate}
 
\begin{figure}[H]
  \centering
  \includegraphics[width=0.72\textwidth]{report/leaderboard.jpeg}
  \caption{Wins leaderboard bar chart as displayed in the pygame window.}
\end{figure}
 

\section{Retrospective}

 
\subsection{Challenges Faced}
 
\textbf{Win detection performance.}
An initial loop-based win checker for Connect 4 was slow on a 7$\times$7 board
with many pieces. We replaced it with
\texttt{np.lib.stride\_tricks.sliding\_window\_view}, which avoids explicit
Python loops entirely and runs in microseconds.
 
\textbf{Diagonal win detection.}
Extracting anti-diagonals from a stride-tricks window was non-obvious.
The solution was to flip the board horizontally with \texttt{np.fliplr} before
applying the same main-diagonal extraction, so both diagonal directions use
identical code paths.
 
\textbf{Transparent PNG graphs over a busy background.}
Early matplotlib figures saved with a white background were jarring against the
retro arcade scene. Setting \texttt{fig.patch.set\_alpha(0)} and
\texttt{ax.set\_facecolor((0,0,0,0))} made the figure areas fully transparent,
and a semi-transparent black pygame surface was rendered behind each graph to
ensure readability.
 
\textbf{Custom pixel font in matplotlib.}
\texttt{matplotlib} does not load \texttt{.ttf} files from the working directory
automatically. We had to call \texttt{fm.fontManager.addfont()} and construct
a \texttt{FontProperties} object, passing it explicitly to every text element.
Forgetting a single element reverted to the default sans-serif font.
 
\textbf{CSV data hygiene.}
Early test runs left rows with ``QUIT'' or ``None'' as winner/loser values.
Both the Python recorder and the AWK script had to be hardened to skip these
rows, otherwise they appeared in the leaderboard as real players.
 
\textbf{Bash password hashing portability.}
\texttt{sha256sum} appends a filename (\texttt{-}) to its output. Forgetting the
\texttt{awk '\{print \$1\}'} filter stored the dash character as part of the
hash, so no login ever succeeded. This was found by inspecting
\texttt{users.tsv} manually.
 
\subsection{What We Would Do Differently}
 
Given ten more hours we would make the following improvements:
 
\begin{itemize}[leftmargin=*]
  \item \textbf{Persist leaderboard state in SQLite} rather than a flat CSV,
        enabling proper JOIN queries and avoiding the manual AWK aggregation.
  \item \textbf{Add an AI opponent} for single-player mode, at minimum a
        random-move bot and ideally a minimax agent for Connect 4.
  \item \textbf{Sound effects and background music} using
        \texttt{pygame.mixer} to reinforce the retro arcade theme.
  \item \textbf{Animated transitions} between screens (fade in/out) to replace
        the current immediate swap.
  \item \textbf{In-game rematch button} so players do not need to navigate
        back through the hub and sort-menu just to play again.
  \item \textbf{Better authentication}: store a salt alongside each hash to
        prevent rainbow-table attacks; use \texttt{bcrypt} via Python instead
        of shell \texttt{sha256sum}.
  \item \textbf{Unit tests} for the game logic (win detection, flip logic,
        valid-move generation) using \texttt{pytest}.
\end{itemize}
 

\section{References}

 
\begin{enumerate}[leftmargin=*, label={[\arabic*]}]
  \item Pygame Documentation. \textit{pygame.org}. \\
        \url{https://www.pygame.org/docs/}
 
  \item Matplotlib Documentation. \textit{matplotlib.org}. \\
        \url{https://matplotlib.org/stable/contents.html}
 
  \item NumPy \texttt{sliding\_window\_view}. \textit{numpy.org}. \\
        \url{https://numpy.org/doc/stable/reference/generated/numpy.lib.stride_tricks.sliding_window_view.html}
 
  \item Othello / Reversi rules. \textit{World Othello Federation}. \\
        \url{https://www.worldothello.org/about/about-othello/othello-rules}
 
  \item SHA-256 in Bash. \textit{GNU Coreutils Manual}. \\
        \url{https://www.gnu.org/software/coreutils/manual/coreutils.html}
 
  \item AWK Programming Language. \textit{The GNU Awk User's Guide}. \\
        \url{https://www.gnu.org/software/gawk/manual/gawk.html}
 
  \item Press Start 2P Font. \textit{Google Fonts}. \\
        \url{https://fonts.google.com/specimen/Press+Start+2P}
 
  \item matplotlib \texttt{FontProperties}. \textit{matplotlib.org}. \\
        \url{https://matplotlib.org/stable/api/font_manager_api.html}
\end{enumerate}
 
\end{document}